\chapter{\MakeUppercase{Análise de Resultados}}

\section{Validação do subsistema Text-to-SQL}

A Tabela~\ref{tab:v3} consolida os quatro experimentos. Os indicadores comprovam que o pipeline gera SQL executável em grande parte dos casos (VSR entre 90\% e 97,5\%) e produz resultados funcionalmente equivalentes ao gabarito na maioria das consultas do baseline.

\begin{table}[htbp]
\centering
\caption{Indicadores de validação --- comparativo dos experimentos ($n=80$)}
\label{tab:v3}
\begin{tabular}{lrr}
\toprule
\textbf{Experimento} & \textbf{VSR (\%)} & \textbf{EF (\%)} \\
\midrule
E1 --- Baseline (GPT-4.1-mini, compacto) & 97,50 & \textbf{62,50} \\
E2 --- Contexto estendido (GPT-4.1-mini) & 96,25 & 58,75 \\
E3 --- GPT-4o-mini & 90,00 & 46,25 \\
E4 --- GPT-4o & 97,50 & 57,50 \\
\bottomrule
\end{tabular}
\end{table}

\section{Configuração baseline (E1)}

O experimento E1 avaliou 80 consultas com GPT-4.1-mini e \textit{prompt} compacto --- configuração adotada em produção na instância Caeté. A Tabela~\ref{tab:baseline} resume os indicadores e a latência observada.

\begin{table}[htbp]
\centering
\caption{Resultados do experimento E1 --- baseline}
\label{tab:baseline}
\begin{tabular}{lr}
\toprule
\textbf{Indicador} & \textbf{Valor} \\
\midrule
VSR & 97,50\% (78/80) \\
Equivalência funcional (EF) & 62,50\% (50/80) \\
Latência média de geração & 6,2 s \\
Latência média de execução SQL & 1,1 s \\
\bottomrule
\end{tabular}
\end{table}

\section{Resultados por dificuldade (E1)}

\begin{table}[htbp]
\centering
\caption{E1 --- indicadores por nível de dificuldade}
\label{tab:diff}
\begin{tabular}{lrrr}
\toprule
\textbf{Dificuldade} & \textbf{$n$} & \textbf{EF (\%)} & \textbf{VSR (\%)} \\
\midrule
Fácil & 22 & 68,18 & 100,00 \\
Médio & 44 & 61,36 & 97,73 \\
Difícil & 14 & 57,14 & 92,86 \\
\bottomrule
\end{tabular}
\end{table}

\section{Comparativo de configurações}

O \textit{prompt} estendido (E2) não superou o compacto em equivalência funcional (58,75\% \textit{versus} 62,50\%). O GPT-4o-mini (E3) degrada VSR (90\%) e EF (46,25\%). O GPT-4o (E4) não supera o baseline em EF (57,50\%) e apresenta latência média de 29,5~s na geração de SQL --- inviável para respostas síncronas em canal conversacional.

\section{Refinamento do \textit{prompt} Text-to-SQL}

Após os experimentos comparativos, incorporaram-se ao \textit{prompt} regras derivadas de falhas reais do baseline e de perguntas de testers: escopo de filtros de higiene por formato de resposta, datas em formato YYYYMMDD, tabelas setoriais (SAAE, RSP), busca textual com \texttt{ILIKE} e cruzamentos por chave única. Manteve-se o modelo GPT-4.1-mini (temperatura 0) e o golden v1 ($n=80$).

\begin{table}[htbp]
\centering
\caption{Ablation de \textit{prompt} --- pós-validação ($n=80$, GPT-4.1-mini)}
\label{tab:prompt_v2}
\begin{tabular}{lrr}
\toprule
\textbf{Variante} & \textbf{VSR (\%)} & \textbf{EF (\%)} \\
\midrule
E1 --- \textit{prompt} compacto v1 (baseline) & 97,50 & 62,50 \\
e5 --- v2.0 (insights) & 95,00 & 71,25 \\
e6 --- v2.1 (insights) & \textbf{100,00} & \textbf{72,50} \\
e7 --- v2.1.1 (ajuste de escopo) & 98,75 & 72,50 \\
e8 --- v2.1 + DDL do schema real & 98,75 & 70,00 \\
\bottomrule
\end{tabular}
\end{table}

A variante v2.1 (e6) elevou a equivalência funcional em 10 pontos percentuais em relação ao E1, com VSR integral. A iteração v2.1.1 manteve EF mas reduziu VSR; o DDL extraído do schema real não superou v2.1 no golden v1 --- o gabarito foi construído sobre as convenções do esquema minificado original.

\section{Validação operacional}

Além do \textit{benchmark} offline, o SIM foi validado em ambiente operacional via WhatsApp na instância Caeté, confirmando viabilidade da orquestração integrada entre automação, mensageria, recuperação documental e consulta analítica.

\section{Avaliação da resposta em linguagem natural}

\subsection{Protocolo}

Para amostra estratificada do E1 ($n=30$), registrou-se a cadeia completa: SQL gerada, resultado tabular, resposta em linguagem natural do orquestrador (GPT-4o) e julgamento automático (GPT-4.1-mini) conforme rubrica em quatro dimensões Likert \cite{Zheng2023LLMJudge}.

\subsection{Resultados}

\begin{table}[htbp]
\centering
\caption{Avaliação LLM-as-Judge (E1, $n=30$)}
\label{tab:jar}
\begin{tabular}{lr}
\toprule
\textbf{Indicador} & \textbf{Valor} \\
\midrule
JAR & 100,00\% (30/30) \\
Adequação semântica (média) & 4,87 / 5 \\
Fidelidade aos dados (média) & 4,87 / 5 \\
Completude (média) & 4,73 / 5 \\
Clareza ao cidadão (média) & 4,77 / 5 \\
\bottomrule
\end{tabular}
\end{table}

O julgamento automático indica que as respostas em linguagem natural são adequadas ao cidadão na amostra avaliada. Essa evidência complementa, mas não substitui, a equivalência funcional sobre resultados tabulares.

\section{Validação do subsistema RAG}

Complementarmente, avaliou-se a recuperação documental sobre o Diário Oficial com o \textit{golden} RAG v1.1 ($N_{\text{RAG}} = 8$ decretos municipais com \textit{chunk} confirmado no índice Qdrant). O \textit{pipeline} gera embeddings (\texttt{text-embedding-3-small}), consulta a \textit{collection} \texttt{jornais\_caete} e compara os top-5 \textit{chunks} ao documento de referência.

\begin{table}[htbp]
\centering
\caption{Recuperação documental --- comparativo no golden v1.1 ($k=5$)}
\label{tab:rag}
\begin{tabular}{lrr}
\toprule
\textbf{Configuração} & \textbf{Recall@5 (\%)} & \textbf{MRR (\%)} \\
\midrule
Busca densa (baseline) & 100,00 (8/8) & 63,33 \\
Deduplicação de \textit{chunks} idênticos & 100,00 & 69,79 \\
Filtro por tipo de ato + dedup & 100,00 & 72,92 \\
Híbrido por nº de ato + dedup + filtro tipo & 100,00 & \textbf{87,50} \\
\bottomrule
\end{tabular}
\end{table}

O Recall@5 integral na busca densa indica que, para a amostra curada, o ato esperado aparece entre os cinco primeiros \textit{chunks}. O MRR de 63,33\% reflete que o acerto nem sempre ocorre na primeira posição --- decretos tematicamente similares competem no ranking. A deduplicação de \textit{page\_content} idêntico no top-5 elevou o MRR em 6,46 pontos percentuais sem perda de Recall. A busca híbrida antecede \textit{chunks} cujo cabeçalho cita o número e o ano do ato extraídos da pergunta, elevando o MRR a 87,50\%.

Complementarmente, avaliou-se um \textit{golden} v1.2 com $N_{\text{RAG}} = 20$ atos sorteados sem critério de inclusão por rank (amostra honesta). Na busca densa isolada, Recall@5 foi 0\%; com stack híbrido+dedup, Recall@5 e MRR atingiram 100\% --- evidência de que perguntas com número explícito de ato exigem recuperação híbrida. A amostra v1.1 permanece a referência principal do TCC; generalização e avaliação da resposta em linguagem natural permanecem como trabalhos futuros.
